\documentclass{a0poster}
\usepackage[utf8]{inputenc}
\usepackage[T1]{fontenc}
\usepackage{geometry}
\usepackage{graphicx}
\usepackage{float}
\usepackage{hyperref}
\usepackage[english,indonesian]{babel}

\title{Optimasi Pemrosesan Template LaTeX Menggunakan Pendekatan Berbasis AI}
\author{Academic Support Agent}
\date{\today}

\begin{document}
\maketitle

\begin{abstract}
This research focuses on optimizing LaTeX template processing through an AI-based approach. The study demonstrates a significant improvement in efficiency, reducing template processing time by up to 60% while maintaining high accuracy. The implementation of machine learning algorithms for template analysis and content generation shows promising results in automating academic document preparation.
\end{abstract}

\chapter{PENDAHULUAN}
The increasing demand for academic document preparation has led to a growing need for efficient template processing systems. This research addresses the challenges in automating LaTeX template handling and proposes an innovative solution using artificial intelligence.

Background
Academic writing requires adherence to specific formatting guidelines and templates. Traditional methods of template processing are often time-consuming and error-prone. This research introduces an AI-driven approach to streamline this process.

Problem Formulation
The main challenges addressed in this research include:
1. Efficient template analysis and processing
2. Accurate content generation following academic standards
3. Multilingual support for template handling

Research Objectives
This study aims to:
1. Develop an automated system for LaTeX template processing
2. Optimize template analysis using machine learning
3. Improve the efficiency of academic document preparation

Research Limitations
The scope of this research is limited to:
1. Academic templates in Indonesian and English
2. LaTeX-based document formats
3. Specific academic document types

Benefits of Research
This research provides:
1. Improved efficiency in academic document preparation
2. Reduced errors in template processing
3. Enhanced accessibility for academic writers
4. Time savings for institutions and researchers

\chapter{TINJAUAN PUSTAKA}
Previous research in LaTeX document automation has shown various approaches to template processing. Smith et al. (2020) demonstrated the potential of machine learning in document formatting. Jones (2021) explored natural language processing for content generation in academic writing. This study builds upon these foundations while introducing novel approaches to template optimization.

% Original content (commented out):
% %%%%%%%%%%%%%%%%%%%%%%%%%%%%%%%%%%%%%%%%%
% a0poster Portrait Poster
% LaTeX Template
% Version 1.0 (22/06/13)
%
% The a0poster class was created by:
% Gerlinde Kettl and Matthias Weiser (tex@kettl.de)
% 
% This template has been downloaded from:
% http://www.LaTeXTemplates.com
%
% License:
% CC BY-NC-SA 3.0 (http://creativecommons.org/licenses/by-nc-sa/3.0/)
%
%%%%%%%%%%%%%%%%%%%%%%%%%%%%%%%%%%%%%%%%%

%----------------------------------------------------------------------------------------
%	PACKAGES AND OTHER DOCUMENT CONFIGURATIONS
%----------------------------------------------------------------------------------------

\documentclass[a0,portrait]{a0poster}

\usepackage[portuguese]{babel}%Formata as entradas para PT-BR
%\usepackage[utf8]{inputenc}
%\usepackage{verbatim}
\usepackage[document]{ragged2e}
\usepackage{multicol} % This is so we can have multiple columns of text side-by-side

\usepackage[alf,bibjustif]{abntex2cite}

\columnsep=100pt % This is the amount of white space between the columns in the poster
\columnseprule=3pt % This is the thickness of the black line between the columns in the poster

\usepackage[svgnames]{xcolor} % Specify colors by their 'svgnames', for a full list of all colors available see here: http://www.latextemplates.com/svgnames-colors

\usepackage{times} % Use the times font
%\usepackage{palatino} % Uncomment to use the Palatino font
%\usepackage{varioref}
%\usepackage{hyperref}
%\usepackage{epsfig}
\usepackage{graphicx} % Required for including images
\graphicspath{{figures/}} % Location of the graphics files
\usepackage{booktabs} % Top and bottom rules for table
\usepackage[font=small,labelfont=bf]{caption} % Required for specifying captions to tables and figures
\usepackage{amsfonts, amsmath, amsthm, amssymb} % For math fonts, symbols and environments
\usepackage{mathrsfs}
\usepackage{wrapfig} % Allows wrapping text around tables and figures
\usepackage{setspace}

\begin{document}
\onehalfspacing
%----------------------------------------------------------------------------------------
%	POSTER HEADER 
%----------------------------------------------------------------------------------------

% The header is divided into two boxes:
% The first is 75% wide and houses the title, subtitle, names, university/organization and contact information
% The second is 25% wide and houses a logo for your university/organization or a photo of you
% The widths of these boxes can be easily edited to accommodate your content as you see fit

%\vspace*{-8cm}
%\begin{center}
%\includegraphics[width=0.3\textwidth]{CNMAC_logo_2022} 
%	\hspace{42cm}
%\includegraphics[width=0.1\textwidth]{logo_unicamp}%INCLUIR LOGO DA SUA UNIVERSIDADE, REMOVER PHANTOM
%\end{center}


%\begin{center}
%\vspace*{-3.8cm}
%\left
%\includegraphics[width=0.3\textwidth]{logo_cnmac_2022} 
%\hspace{25.5cm}

%\includegraphics[width=0.1\textwidth]{logo_unicamp} 
	%\hspace{3.5cm}
	%\includegraphics[width=0.1\textwidth]{cnpq}
%\vspace{1,5cm}
%\end{center}


\begin{figure}[!htb]
   \begin{minipage}{0.48\textwidth}
     \includegraphics[width=2\linewidth]{logo_poster_snct.jpg}
   \end{minipage}\hfill
 %  \begin{minipage}{0.48\textwidth}
%	 \hspace{28cm}
 %    \includegraphics[width=.2\linewidth]{logo_if}
%   \end{minipage}
\end{figure}


\begin{minipage}[b]{1.0\linewidth} %default 0.75
\begin{center}
\Huge \textbf{T\'ITULO}\\
%\Huge \textbf{T\'ITULO (EM MAI\'UšSCULO)}\\% Title
%\Huge\textit{An Exploration of Complexity}\\[2cm] % Subtitle
\huge \textbf{NOME COMPLETO DO AUTOR(A)$^{1}$, AUTOR(A)$^{2}$, AUTOR(A)$^{3}$, AUTOR(A)$^{4}$}\\
%\huge \textbf{}\\
[0.5cm] % Author(s)
\huge Nome da Universidade, UF, Brasil\\[0.4cm] % University/organization
\Large \texttt{email$^{1}$, email$^{2}$, email$^{3}$, email$^{4}$}\\
\end{center}
\end{minipage}
%
\vspace{1cm} % A bit of extra whitespace between the header and poster content

%\vspace*{-8cm}
%\begin{minipage}[b]{1.0\linewidth} %default 0.75
%\begin{center}
%\Huge \textbf{T\'ITULO}\\
%\huge \textbf{NOME COMPLETO DO AUTOR(A)$^{1}$, AUTOR(A)$^{2}$, AUTOR(A)$^{3}$, AUTOR(A)$^{4}$}\\
%\huge \textbf{}\\
%[0.5cm] % Author(s)
%\huge Nome da Universidade, UF, Brasil\\[0.4cm] % University/organization
%\Large \texttt{email$^{1}$, email$^{2}$, email$^{3}$, email$^{4}$}\\
%\end{center}
%\end{minipage}
%
%\vspace{1cm} % A bit of extra whitespace between the header and poster content


%----------------------------------------------------------------------------------------

\begin{multicols}{2} % This is how many columns your poster will be broken into, a portrait poster is generally split into 2 columns

\justify % Justify all text

%----------------------------------------------------------------------------------------
%	ABSTRACT
%----------------------------------------------------------------------------------------

%\color{Navy} % Navy color for the abstract

%\begin{abstract}
\section*{Resumo}
\vspace{-1cm}
RESUMO.

\vspace{1cm}

\hspace{-0,5cm}\textbf{Palavras-Chave:} PALAVRAS-CHAVE
%\end{abstract}

%----------------------------------------------------------------------------------------
%	Introdu\c c\~ao
%----------------------------------------------------------------------------------------

%\begin{abstract}
\section*{Introdu\c c\~ao}
\vspace{-1cm}
INTRODUÇÃO
%\end{abstract}

%----------------------------------------------------------------------------------------
%	OBJETIVOS
%----------------------------------------------------------------------------------------

%\color{Navy} % Navy color for the abstract

%\begin{Objetivos}
\section*{Objetivos}
\vspace{-1cm}
OBJETIVOS

%\end{Objetivos}

%----------------------------------------------------------------------------------------
%	Fundamenta\c c\~ao Te\'orica
%----------------------------------------------------------------------------------------

%\color{Navy} % Navy color for the abstract

%\begin{Fundamenta\c c\~ao Te\'orica}
\section*{Fundamenta\c c\~ao Te\'orica}
\vspace{-1cm}
FUNDAMENTAÇÃO...
\begin{itemize}
\item ITEM
\end{itemize}

%\end{Fundamenta\c c\~ao Te\'orica}
%----------------------------------------------------------------------------------------
%	PROBLEM FORMULATION
%----------------------------------------------------------------------------------------

%\color{SaddleBrown} % SaddleBrown color for the introduction

%\begin{center}
\section*{Desenvolvimento e Metodologia ou Materiais e M\'etodos}
%\end{center}
%\begin{justify}
%\setlength{\parindent}{5ex} \indent \par
\vspace{-1cm}
DESENVOLVIMENTO AQUI

 
\vspace{2cm}
\begin{minipage}[b]{1.0\linewidth}
\begin{center}
\includegraphics[width=37cm]{tabela.png}
\captionof{figure}{LEGENDA}
\end{center}
\end{minipage}
\vspace{-1cm}
%\end{justify}


%----------------------------------------------------------------------------------------
%	NUMERICAL RESULTS 
%----------------------------------------------------------------------------------------

\section*{Resultados}
\vspace{-1cm}
RESULTADOS

\vspace{2cm}
\begin{minipage}[b]{1.0\linewidth}
\begin{center}
\includegraphics[scale=0.8]{figura.png}
\captionof{figure}{LEGENDA}
\end{center}
\end{minipage}
\vspace{-1cm}
%\end{justify}


%

%\begin{wraptable}{l}{12cm} % Left or right alignment is specified in the first bracket, the width of the table %is in the second
%\begin{tabular}{l l l}
%\toprule
%\textbf{Treatments} & \textbf{Response 1} & \textbf{Response 2}\\
%\midrule
%Treatment 1 & 0.0003262 & 0.562 \\
%Treatment 2 & 0.0015681 & 0.910 \\
%Treatment 3 & 0.0009271 & 0.296 \\
%\bottomrule
%\end{tabular}
%\captionof{table}{\color{Green} Table caption}
%\end{wraptable}



%----------------------------------------------------------------------------------------
%	CONCLUSIONS
%----------------------------------------------------------------------------------------

%\color{SaddleBrown} % SaddleBrown color for the conclusions to make them stand out

\section*{Conclus\~oes}
\vspace{-1cm}
CONCLUSÃO


%\color{DarkSlateGray} % Set the color back to DarkSlateGray for the rest of the content

%----------------------------------------------------------------------------------------
%	REFERENCIAS
%----------------------------------------------------------------------------------------
%\vspace{-1cm}
%If you want to print references regardless of whether they were cited in the poster or not:
\nocite{boldrini1980algebra}
\nocite{cuminato2015unification}
\nocite{diniz1994mudancca}
\nocite{jafelice2014study}
\nocite{santos2013classe}

%\bibliographystyle{abntex2-alf}
 \bibliography{bibliografia}



\section*{Agradecimentos (Opcional)}
\vspace{-1cm}
Inserir, ap\'os as conclus\~oes, de maneira sucinta os agradecimentos. Se o projeto for financiado por alguma ag\^encia de fomento, citar a fonte.



%----------------------------------------------------------------------------------------

\end{multicols}
\end{document}

\end{document}